\begin{theorem}[S-36: Fractional determinants do not repair the balanced-class denominator]
\label{thm:semireg-rm-exact-denominator}
Let $(X,\theta)$ be a principally polarized complex
abelian fourfold with Rosati-invariant real multiplication
by a genuine real quadratic field $F$ of type $(2,2)$.
Assume that its full ring of integers $\mathcal O_F$
acts on $H_1(X,\mathbf Z)$ and that
$\operatorname{NS}(X)_{\mathbf Q}=F\theta$.
Let $\Delta_F$ be the positive field discriminant and set
\[
 d_F=\begin{cases}
       \Delta_F,&\Delta_F\text{ odd},\\
       \Delta_F/2,&\Delta_F\text{ even}.
      \end{cases}
 \qquad C=\frac{\theta_1^2+\theta_2^2}{2}.
 \tag{S36.1}
\]
Suppose a positive-rank perfect object $E$ satisfies
\[
 e^{-c_1(E)/N}\operatorname{ch}(E)
       =N\left(1-qC+q^2\frac{\theta^4}{24}\right),
 \qquad N=\operatorname{rk}(E)>0,
 \qquad q\in\mathbf Q_{>0},
 \qquad u=Nq.
 \tag{S36.2}
\]
Then
\[
 u\in d_F\mathbf Z_{>0},\qquad uC\in H^4(X,\mathbf Z).
 \tag{S36.3}
\]
The integer $d_F$ is exactly the least positive integer
which makes $d_FC$ integral. Thus a fractional determinant
correction cannot cancel any part of this denominator.
No simplicity or semiregularity is assumed in (S36.3).

There is also a determinant restriction. Define
\[
 s(N)=\prod_p p^{\lceil v_p(N)/2\rceil},\qquad
 c(N)=\prod_p p^{\lfloor v_p(N)/2\rfloor},
 \qquad N=s(N)c(N).
\]
The integral divisor class $c_1(E)$ is divisible by
$s(N)$ in $\operatorname{NS}(X)$, so the denominator of
$-c_1(E)/N$ divides $c(N)$. In particular, that normalizing
class is integral whenever $N$ is squarefree.

If $E$ has a simple gluable fully semiregular lift for an
effective finite coherent target action $G$, S-22 gives
\[
 |G|\in\{u^2,2u^2\},\qquad d_F^2\mid |G|.
 \tag{S36.4}
\]
For $F=\mathbf Q(\sqrt2)$ this requires $4\mid u$;
the group order is a multiple of $16$ in the Euler-four
case and a multiple of $32$ in the Euler-two case.
These are necessary conditions, not existence assertions.
\end{theorem}
\begin{proof}
Write $\mathcal O_F=\mathbf Z[\omega]$, put
$T=\operatorname{Tr}_{F/\mathbf Q}(\omega)$ and
$U=\operatorname{Nm}_{F/\mathbf Q}(\omega)$, and set
$\delta=2\omega-T$, so $\delta^2=\Delta_F$.
Use the standard integral basis, for which $T=1$ if
$\Delta_F$ is odd and $T=0$ otherwise.
The trace pairing $(a,b)\mapsto\operatorname{Tr}(ab/\delta)$
is perfect on $\mathcal O_F$: in the basis $1,\omega$
its matrix is
\[
 \begin{pmatrix}0&1\\1&T\end{pmatrix}.
 \tag{S36.5}
\]
This explicit matrix avoids any choice of a global free
basis for the rank-four $\mathcal O_F$-lattice.

Fix a rational prime $p$ and put
$R=\mathcal O_F\otimes\mathbf Z_p$.
This is a DVR or a product of two DVRs, and
$H_1(X,\mathbf Z_p)$ is free of rank four over each
corresponding factor. Rosati symmetry and (S36.5)
uniquely express the principal alternating form as
\[
 \theta=\operatorname{Tr}_{F\otimes\mathbf Q_p/\mathbf Q_p}
                       (\psi/\delta),
 \tag{S36.6}
\]
where $\psi$ is a perfect alternating $R$-bilinear form.
The trace argument and dual-lattice equality are the
same as in S-32 and S-33; alternation holds also at two,
since the ground fields have characteristic zero.
Symplectic elimination over the local factors gives
two orthogonal $R$-hyperbolic planes.

Put $\ell=c_1(E)=\theta_r$ for $r\in F$.
Integrality of $\ell$ and unimodularity of $\theta$
imply that $r$ preserves the integral lattice. In the
local free $R$-module this forces $r\in R$; hence
$r\in\mathcal O_F$ globally.
For $x\in R$ define $f(x)=\operatorname{Tr}(rx/\delta)$.
On a mixed four-tuple $(xe_1,f_1,ye_2,f_2)$ from the
two hyperbolic planes one has
\[
 C(xe_1,f_1,ye_2,f_2)=\frac{\operatorname{Tr}(xy)}{\Delta_F},
 \qquad
 \frac{\ell^2}{2}(xe_1,f_1,ye_2,f_2)=f(x)f(y).
 \tag{S36.7}
\]
The matrix of the first form on $1,\omega$ is
\[
 A=\frac1{\Delta_F}
       \begin{pmatrix}2&T\\T&T^2-2U\end{pmatrix},
 \qquad \det A=\frac1{\Delta_F}.
 \tag{S36.8}
\]
The full four-form $C$ has only the two-plane mixed
components: on each hyperbolic plane each embedding
component of $\theta$ has square zero. Its other
mixed evaluations are obtained by allowing the
$R$-valued alternating products in (S36.7) to range
over $R$. Thus integrality of a scalar multiple of
$C$ is equivalent locally to integrality of the same
scalar multiple of $A$.

The second Chern identity following from (S36.2) says
\[
 uC-\frac{\ell^2}{2N}\in H^4(X,\mathbf Z),
 \tag{S36.9}
\]
because $\ell^2/2$ is integral on a torus. Consequently,
with $z=(f(1),f(\omega))^t\in\mathbf Z_p^2$,
\[
 uA-\frac{zz^t}{N}\in\operatorname{Mat}_2(\mathbf Z_p).
 \tag{S36.10}
\]
We determine the possible valuation of $u$ at every prime.

If $p\nmid\Delta_F$, the matrix $A$ is integral and
unimodular. If $v_p(u)<0$, multiply (S36.10) by
$p^{-v_p(u)}$ and reduce modulo $p$. The rank-one
matrix on the right would become a nonzero scalar
multiple of the rank-two reduction of $A$, which is
impossible. Thus $v_p(u)\ge0$ at every unramified prime,
including an unramified prime two.

At an odd ramified prime write $F\otimes\mathbf Q_p
=\mathbf Q_p(\sqrt d)$, where $d$ is the squarefree
radicand and $v_p(d)=1$. The integral basis
$1,\sqrt d$ gives, after the same hyperbolic
normalization, the three congruences
\[
 \frac{b^2}{N}\equiv\frac{u}{2d},\qquad
 \frac{a^2}{N}\equiv\frac u2,\qquad
 \frac{ab}{N}\equiv0\pmod{\mathbf Z_p},
 \qquad r=a+b\sqrt d.
 \tag{S36.11}
\]
If $v_p(u)<0$, the first two fractions have valuations
$v_p(u)-1$ and $v_p(u)$, respectively. Their difference
is odd, whereas the difference of the valuations of
$b^2/N$ and $a^2/N$ is even. This is impossible.
The remaining value $v_p(u)=0$ is excluded by S-33's
Euler-integrality argument. Therefore $v_p(u)\ge1$.

Suppose next that two ramifies, so
$\mathcal O_F\otimes\mathbf Z_2=\mathbf Z_2[\sqrt d]$
with squarefree $d\equiv2$ or $3\pmod4$.
Congruences (S36.11) hold at two as well.
If $d$ is odd and $v_2(u)\le0$, the first two fractions
have the same negative valuation $v_2(u)-1$.
Their square roots then force
$v_2(ab/N)=v_2(u)-1<0$, contrary to the third congruence.
Hence $v_2(u)\ge1$ when $d\equiv3\pmod4$.

If $d$ is even and $v_2(u)\le0$, the first two
negative valuations differ by one, contradicting the
same parity argument. It remains to exclude
$v_2(u)=1$. Put $m=v_2(N)$. The first congruence
forces
\[
 m=2v_2(b)+1>0,
 \qquad v_2(a)\ge(m+1)/2,
 \tag{S36.12}
\]
where the second assertion follows from $a^2/N$ being
integral. It follows that $r^2/N$ is a unit in the
ramified quadratic valuation ring: $db^2/N$ is a unit,
while $a^2/N$ and $2ab\sqrt d/N$ are nonunits.
Since $u$ is even, $u-r^2/N$ is also a unit.
The exact Riemann--Roch identity is
\[
 \chi(E)=\frac{\operatorname{Nm}_{F/\mathbf Q}(Nu-r^2)}{N^3}.
 \tag{S36.13}
\]
Indeed, expansion of $e^{\ell/N}$ in (S36.2) uses
$\ell^4/24=\operatorname{Nm}(r)^2\theta^4/24$ and
$C\ell^2/2=\operatorname{Tr}(r^2)\theta^4/24$.
The numerator in (S36.13) has two-adic valuation $2m$
because $Nu-r^2=N(u-r^2/N)$. Thus $v_2(\chi(E))=-m<0$,
contradicting the integer Euler characteristic.
We conclude that $v_2(u)\ge2$ for even $d$.

These local conclusions say exactly that $u$ is a
positive integer divisible by $d_F$. Matrix (S36.8)
also shows that this is the exact denominator of $C$:
the common divisor of its integral numerators is
$\gcd(2,T)$, which is one for odd discriminant and two
for even discriminant. The mixed-component description
following (S36.8) proves both necessity and sufficiency
for integral scalar multiples of $C$. This establishes
(S36.3).

Now $uC$ is integral, so (S36.9) implies that
$\ell^2/(2N)$ is integral. Formula (S36.7) then shows
that every entry of $zz^t/N$ is integral at $p$.
In particular,
$v_p(f(1)),v_p(f(\omega))\ge\lceil v_p(N)/2\rceil$.
The trace-dual matrix (S36.5) is unimodular, hence
$r\in p^{\lceil v_p(N)/2\rceil}R$.
Doing this at every prime gives $r\in s(N)\mathcal O_F$.
The associated divisor class $\theta_{r/s(N)}$ is
integral, proving the determinant and normalizing-class
claims. Finally S-22 gives $4u^2/|G|\in\{2,4\}$,
which is exactly (S36.4).
\end{proof}
